\section{From certificates to experiments}\label{sec:diagnostic}
The scalar maximum in \eqref{eq:program} determines whether to continue. Its optimizer determines what to observe. We derive the sampling rule from the same variables that implement the dual norm in the cone program.

\subsection{Implementing the maximizer}
For a single core action string the observation outcomes are disjoint. A signed coefficient on those outcomes can therefore be measured simply by executing the string and reporting a lookup-table score. An adaptive continuation creates a second possibility. A policy can prescribe one string after one observation and a different string after another. To sample such a policy it suffices to know how much action probability reaches each node of its test tree. The key convex step is that a coefficient in the dual ball can be split between these two kinds of experiments.

\begin{lemma}[A dual coefficient is an experiment]\label{lem:score}
At any stage, $\dual c\le b$ has a polynomial-size conic representation. A feasible representation yields an executable controller $K_c$, a mass $M_c\le\sqrt{U_h}\,b$ and an observable score $|f_c|\le1$ such that, for every valid state $q$,
\begin{equation}
 c^\top q=M_c\,\E_q^{K_c}f_c .
 \label{eq:scored}
\end{equation}
\end{lemma}
\begin{proof}
Suppress the stage index. Split $c=s+t$ with cost $b_1+b_2\le b$. For the group part impose $|s_j|\le d_{\alpha(j)}$, $d\ge0$, $\|d\|_2\le b_1$. For the second part place a nonnegative realization weight $p$ on the test-prefix tree, with root mass $b_2$. At an observation branch copy the incoming weight to each possible observation. At an action decision split it over actions:
\[
 p(\ell o)=p(\ell),\qquad
 \textstyle\sum_a p(\ell oa)=p(\ell o).
\]
Include each action at a decision node, allowing exits from the tree. Require $|t_j|\le p(j)$ on maximal leaves, and $t_j=0$ elsewhere. At a positive-mass decision the ratio $p(\ell oa)/p(\ell o)$ implements its policy, so $p(j)=b_2\pi(j)$.

To see exactness, the group body's support, for unit cost, is $N_{12}$ by groupwise maximization and Cauchy--Schwarz. The signed flow body's support is $N_\Pi$: choose a maximizing policy and leaf signs. Splitting a unit cost between the two bodies forms their convex hull, whose support is the maximum of the two support functions. By \eqref{eq:fused} this is precisely the fused-dual unit ball. There are $O(AHD)$ flow variables per coefficient block. The remaining history constraint in \eqref{eq:program} is one Euclidean norm bound on a vector of linear evaluations.

For the group experiment take $M_s=\sum_\alpha d_\alpha$. Draw string $\alpha$ with probability $d_\alpha/M_s$, apply it and return $s_j/d_\alpha$ on listed outcome $j$, or zero on an unlisted outcome. Its expected score times $M_s$ is $s^\top q$. For the flow experiment run the realization policy. On reaching a maximal leaf $j$, return $t_j/p(j)$, and return zero on exit. The probability of leaf $j$ is $\pi(j)q_j$, so this experiment has mean $t^\top q/b_2$. Mix the experiments in proportions $M_s:b_2$. Their total mass satisfies $M_s+b_2\le\sqrt{U_h}\,b_1+b_2$. Omit zero-weight branches and use the constant zero score when the total mass vanishes.
\end{proof}

\paragraph{Two diagnostic continuations.}
At the current prefix, fix a feasible optimizer $(u,w)$ and its conic decomposition by a deterministic evaluation rule. Let $M_0=W/2$. To measure the first term in $g_v$, run $K_u$ with score $M_u f_u$. For the second term, observe $o$, then choose an action by a distribution $\alpha(a\mid o,x)\ge M_{w_{oa}}/M_0$. Such a distribution exists by \eqref{eq:geometry} and Lemma~\ref{lem:score}. Put spare probability on a fixed action. Run $K_{w_{oa}}$ and report
\[
 Z_w=\frac{M_{w_{oa}}}{\alpha(a\mid o,x)}\,f_{w_{oa}} .
\]
A zero-mass term has zero score. Each diagnostic score lies in $[-M_0,M_0]$. By \eqref{eq:update}, its conditional mean is
\[
 \E_\theta[Z_w\mid x]=\sum_{o,a} w_{oa}^{\top}z_h^\theta(o,a;x).
\] The same frozen rule measures the fitted and the true values of the linear functional.

Write $D_h^u,D_h^w$ for following $\pi$ up to stage $h$ and then using the indicated diagnostic. Define a logged mixture
\begin{equation}
 \rho=\tfrac13\pi+\tfrac{1}{3H}\sum_{h=1}^H(D_h^u+D_h^w).
 \label{eq:rho}
\end{equation}
Each diagnostic follows $\pi$ to its selected stage, so its complete log records both the history reached and the outcome of the selected test. The full-policy component also records the unmodified trajectory.

\subsection{Transferring the fitted bonus to the real system}
Fix the data and a policy $\pi$. Even with identical physical core coordinates, replacing the fit by the truth changes two mathematical objects. The sampling measure on prefixes changes. At any prefix reached by both models, the values of the lifted coordinates can also change. There is a third issue specific to a data-dependent support: its historical constraint was imposed with fitted feature rows. We first transport that constraint on the stored prefixes. Then we bound the discrepancy of the selected scores on the current policy's prefixes.

\paragraph{Historical feasibility in true coordinates.}
To compare support functions, use the same stored prefixes in the true representation. Let $\cE=2LS+16Y\beta$. Define $\mathcal V_{\st,k}$ by $|v_0|\le W$, $v_h\in\K_h$, and the single constraint
\[
 N_0v_0^2+\sum_{h,i}g_{v_h}(\phi_h^\st(x_i))^2\le\cE .
\]
Let $b_{\st,k}(\tau)$ be the support of this set at the concatenated true features: maximize $v_0+\sum_h g_{v_h}(\phi_h^\st(\tau_{h-1}))$. From \eqref{eq:drift} every tuple feasible in the separate fitted programs belongs to this true set. Indeed its total fitted energy is at most $LS$ and its total squared drift is at most $8Y\beta$ by Lemma~\ref{lem:mle}. Apply $(a+b)^2\le2a^2+2b^2$. This inclusion is uniform over tuples, so it holds when each stage optimizer is selected at the queried prefix.

\paragraph{On-policy and intervened observations.}
If the planner always picks mode $b$ at the right station, trajectory data from that choice have no power to detect the wrong success probability for mode $a$ in \sect{sec:example}. The transition block of $v$ can nevertheless put a large weight on $a$. Running the $w$ diagnostic samples just such a counterfactual action. Each diagnostic log also records the roll-in under $\pi$, so the law of the complete experiment reflects changes in the probability of reaching its starting history.

\begin{proposition}[Diagnostic transfer]\label{prop:transfer}
For the attained numerical reward $\widetilde b_k$, any frozen $\pi$, and its controller \eqref{eq:rho}, on the calibration event,
\begin{equation}
 \E_{\wh\theta_k,\pi}\widetilde b_k
 \le \E_{\st,\pi} b_{\st,k}
       +3HW\,\TV(P_{\wh\theta_k}^\rho,P_\st^\rho).
 \label{eq:transfer}
\end{equation}
\end{proposition}
\begin{proof}
Write $t_\nu=\TV(P_{\wh\theta_k}^\nu,P_\st^\nu)$ for the distance between complete logged experiments. Freeze the deterministic rule $x\mapsto v_h(x)$ that attains each numerical lower value, including its diagnostic controller and score. For $s\in\{u,w\}$ let $Z_h^s$ be the score returned by $D_h^s$; it lies in $[-W/2,W/2]$. The defining conditional means of these scores, followed by averaging over the roll-in, give
\[
 \E_{\wh\theta_k,\pi}\widetilde b_k
 \le v_0+\sum_{h=1}^H\sum_{s\in\{u,w\}}
      \E_{\wh\theta_k,D_h^s} Z_h^s.
\]
At a true trajectory, the tuple chosen at its visited prefixes is feasible in $\mathcal V_{\st,k}$. Its value is therefore at most $b_{\st,k}$ pointwise, and hence
\[
 v_0+\sum_{h=1}^H\sum_{s\in\{u,w\}}
      \E_{\st,D_h^s} Z_h^s
 \le \E_{\st,\pi}b_{\st,k}.
\]
Because the same frozen score is evaluated on each pair of logged experiment laws and its range has width $W$, the difference of its expectations is at most $Wt_{D_h^s}$. Thus the transfer cost is at most $WS_D$, where $S_D=\sum_h(t_{D_h^u}+t_{D_h^w})$. The mixture label in \eqref{eq:rho} is observed, so $3t_\rho=t_\pi+S_D/H$. Consequently $WS_D\le3HWt_\rho$, as claimed.
\end{proof}

Put $\tau_0=\eps/8$ and $\eta=\tau_0/(6HW)$. A nonstopping round has fitted bonus value above $\tau_0$ with the stated simulation accuracy. Proposition~\ref{prop:transfer} then implies
\begin{equation}
 \E_{\st,\pi^k}b_{\st,k}\ge\tau_0/2
 \quad\text{or}\quad
 \TV(P_{\wh\theta_k}^{\rho_k},P_\st^{\rho_k})\ge\eta.
 \label{eq:either}
\end{equation}
The algorithm collects both batches. The proof can therefore charge the round to whichever inequality holds.
